% q0007.tex — Excuse Me, I Was Here First! (Crowding Out)
\question{
  id        = {0007},
  title     = {Excuse Me, I Was Here First!},
  source    = {OBECON},
  year      = {2026},
  round     = {Seletiva IEO -- Phase A},
  language  = {English},
  topic     = {Macro},
  subtopic  = {Loanable Funds / Crowding Out},
  type      = {objective},
  statement = {%
    In the loanable funds market, suppose the supply and demand for loanable
    funds are given by:
    \[
      S(r) = 400 + 10r \qquad I(r) = 700 - 20r
    \]
    where $r$ is the real interest rate (in percentage points) and $I$ is
    private investment (in billions of dollars). Initially, government
    borrowing is zero ($G = 0$), and the economy is at full employment. The
    government decides to increase spending by \$150 billion, financing this
    entirely through borrowing ($G = 150$).

    By how much does private investment change due to the crowding-out effect?
  },
  choices   = {%
    \item Private investment increases by \$150 billion.
    \item Private investment increases by \$100 billion.
    \item Private investment remains unchanged.
    \item Private investment decreases by \$100 billion.
    \item Private investment decreases by \$150 billion.
  },
  answer    = {D},
  solution  = {%
    \textbf{Step 1 — Initial equilibrium} ($G=0$):
    \[ 400 + 10r = 700 - 20r \implies 30r = 300 \implies r^* = 10\% \]
    \[ I_0 = 700 - 20(10) = 500 \text{ billion} \]

    \textbf{Step 2 — After government borrowing} ($G=150$):
    Total demand for loanable funds $= I(r) + G = (700-20r) + 150 = 850-20r$.
    \[ 400 + 10r = 850 - 20r \implies 30r = 450 \implies r^{**} = 15\% \]
    \[ I_1 = 700 - 20(15) = 400 \text{ billion} \]

    \textbf{Change in private investment:}
    \[ \Delta I = 400 - 500 = -\$100 \text{ billion} \]

    Private investment \textbf{decreases by \$100 billion}. Note that full
    crowding-out would require $\Delta I = -\$150B$; here it is only
    partial because the supply of loanable funds is upward-sloping.
    Answer: \textbf{(D)}.
  },
}
